%% This document is created by 
%%  Dr. Putu Harry Gunawan
%% Template untuk Proposal TA 1 dan TA
%% Template ini digunakan untuk penulisan proposal TA 1 atau TA Fakultas Informatika, Telkom University.
%%  New Update 28 Pebruari 2017 : Bookmarks pdf, Pembimbing hanya satu

\documentclass[a4paper,12pt,oneside]{book}
\usepackage[utf8]{inputenc}
\usepackage{sectsty}
\usepackage{graphicx}
\usepackage{epstopdf}
\usepackage{algorithm}
\usepackage{algpseudocode}
\usepackage{array}
\usepackage[table]{xcolor}
\usepackage{anysize}
\usepackage{amsmath}
\usepackage{amssymb}
\usepackage[bahasa]{babel}
\usepackage{indentfirst} %Spasi untuk paragraf pertama
\usepackage{geometry}
\usepackage{multirow}% http://ctan.org/pkg/multirow
\usepackage{hhline}% http://ctan.org/pkg/hhline
\marginsize{4cm}{3cm}{3cm}{3cm} %{left}{right}{top}{bottom}
\usepackage[compact]{titlesec} 
\usepackage{etoolbox}
\usepackage{natbib} %for Harvard syle
\usepackage[bookmarks,hypertexnames=false,debug]{hyperref}
%\usepackage[pageanchor]{hyperref}
\usepackage{bookmark}


\makeatletter
\patchcmd{\ttlh@hang}{\parindent\z@}{\parindent\z@\leavevmode}{}{}
\patchcmd{\ttlh@hang}{\noindent}{}{}{}
\makeatother

\chapterfont{\centering}
\newcommand{\bigsize}{\fontsize{16pt}{14pt}\selectfont}
\chapterfont{\centering\bigsize\bfseries}
\sectionfont{\large\bfseries}
\usepackage{tikz}
\usetikzlibrary{shapes.geometric, arrows}
%\renewcommand{\chaptertitle}{BAB}
\renewcommand{\thechapter}{\Roman{chapter}}
\renewcommand\thesection{\arabic{chapter}.\arabic{section}}
\renewcommand\thesubsection{\thesection.\arabic{subsection}}
\renewcommand{\theequation}{\arabic{chapter}.\arabic{equation}}
\renewcommand{\thefigure}{\arabic{chapter}.\arabic{figure}}
\renewcommand{\thetable}{\arabic{chapter}.\arabic{table}}

\renewcommand\bibname{Daftar Pustaka}
\addto{\captions bahasa}{\renewcommand{\bibname}{Daftar Pustaka}}
\usepackage{fancyhdr}
\pagestyle{fancy}
\lhead{}
\chead{}
\rhead{}
\lfoot{}
\cfoot{\thepage}
\rfoot{}
\renewcommand{\headrulewidth}{0pt}

\makeatletter

%%%%%%%%%%%%%%%%%%%%%%%%%%%%%%%%%%%%%%%%%%%%%%%%%%%%%%%%%%%%
%
%  Berikut adalah data-data yang wajib diisi oleh mahasiswa
%
%%%%%%%%%%%%%%%%%%%%%%%%%%%%%%%%%%%%%%%%%%%%%%%%%%%%%%%%%%%%

\title{Klasifikasi Data Percakapan dengan Algoritma Naive Bayes untuk Chatbot Youme}\let\Title\@title   %Judul dalam bahasa Indonesia

\newcommand{\EngTitle}{Dialogue Data Classification using Naive Bayes Algorithm for Chatbot Youme}  %Judul dalam bahasa Inggris

\author{Muhammad Husain}  \let\Author\@author  %Nama mhs
\newcommand{\NIM}{1301153626}
\newcommand{\Prodi}{Teknik Informatika}
\newcommand{\KK}{ICM} %UNTUK TA
\newcommand{\Gelar}{Sains Komputasi} % UNTUK TA
\date{2016}           \let\Date\@date %Masukkan hanya tahun saja
\newcommand{\Tanggal}{22} % Tanggal Pengesahan
\newcommand{\Bulan}{Agustus} % Bulan Pengesahan
\newcommand{\PembimbingSatu}{Dr. Calon Pembimbing 1, M.Kom}
\newcommand{\NIPPembimbingSatu}{123456}
\newcommand{\PembimbingDua}{Dr. Calon Pembimbing 2, M.Kom}
\newcommand{\NIPPembimbingDua}{123456}
\newcommand{\Kaprodi}{Dr. Ketua Program Studi}
\newcommand{\NIPKaprodi}{95650581-1}
\newif\ifPembimbingHanyaSatu

%%%% WARNING kode berikut ini diaktifkan Jika pembimbing hanya satu %%%%

%\PembimbingHanyaSatutrue

%%%%%%%%%%%%%%%%%%%%%%%%%%%%%%%%%%%%%%%%%%%%%%%%%%%%%%%%%%%%%%%%%%

\makeatother
\linespread{1}


\begin{document}
%%%%%%%%%%%%%%%%%%%%%%%%%%%%%%%%%%%%%%%%%%%%%%%%%%%%%%%%%%
% Dimulai dengan cover, lembar persetujuan, Abstrak, dll
%%%%%%%%%%%%%%%%%%%%%%%%%%%%%%%%%%%%%%%%%%%%%%%%%%%%%%%%%%%
\pagenumbering{roman} 
\begin{titlepage}
\thispagestyle{empty}
%\vspace*{0.7cm}
\pdfbookmark{Cover}{ }
{\centering
\large
{\bigsize\bf \Title}\\
\vspace{ 2cm}
\rm
\textbf{Proposal Tugas Akhir}\\
\vspace{0.5 cm}
\textbf{Kelas TA 1}\\
\vspace{0.5 cm}
\textbf{\Author}\\ \textbf{NIM: \NIM}\\ 

\vspace{1.5 cm}

\begin{figure}[h]
{\centering {\includegraphics[scale=0.17]{Tel-U-Logo}}\par}
\end{figure}

\vspace{2 cm}
{\bigsize\textbf{Program Studi Sarjana \Prodi}\\
\vspace{0.5 cm}
\textbf{Fakultas Informatika}\\
\vspace{0.5 cm}
\textbf{Universitas Telkom}\\
\vspace{0.5 cm}
\textbf{Bandung}\\
\vspace{0.5 cm}
\textbf{\Date}\\}
}
\pagebreak
\thispagestyle{empty}
\pdfbookmark{Lembar Persetujuan}{ }
{\centering
\textbf{\large Lembar Persetujuan}\\
\vspace{0.5cm}
\textbf{\Title}\\
\vspace{0.5cm}
\textbf{\textit{\EngTitle}}\\
\vspace{0.5cm}
\textbf{\Author}\\
\textbf{NIM: \NIM}\\
\vspace{1cm}


{ Proposal ini diajukan sebagai usulan pembuatan tugas akhir pada\\ Program Studi Sarjana \Prodi\\ Fakultas Informatika Universitas Telkom}\\

\vspace{0.5cm}

{Bandung, \Tanggal\quad \Bulan \quad \Date}\\
{Menyetujui}\\

\vspace{0.5cm}
      \ifPembimbingHanyaSatu
             \begin{center}
            Calon Pembimbing 1
            \end{center}
       \else
          \begin{center}
              \begin{tabular}{  m{8cm}  m{8cm} }
              Calon Pembimbing 1 & Calon Pembimbing 2
              \end{tabular}
           \end{center}
      \fi
\begin{center}
\vspace{2cm}
\ifPembimbingHanyaSatu
     \underline{\PembimbingSatu} \\ 
      NIP: \NIPPembimbingSatu
\else
\begin{tabular}{  m{8cm}  m{8cm} }
\underline{\PembimbingSatu} & \underline{\PembimbingDua} \\ 
NIP: \NIPPembimbingSatu & NIP: \NIPPembimbingDua
\end{tabular}
\fi
\end{center}
\vspace{0.5cm}

}
\pagebreak
\end{titlepage}
\phantomsection
\addcontentsline{toc}{chapter}{Abstrak}
%\pdfbookmark{\abstractname}{Abstrak Indonesia}
\chapter*{Abstract}

 Chatbot is a growing interest currently as a dialog or chatting partner. As a chatting partner, Chatbot are often designed to imitate how human would behave as a chatting partner.Therefore Chatbots are usually used in dialog system for various purpose like Customer Service, Information Acquisition, or Question answering. To be able to fulfill these purpose, correctly defining a topic    of a chat is required
  
\vspace{0.5 cm}
\begin{flushleft}
{\textbf{Kata Kunci:} Shallow, water, equations.}
\end{flushleft}

\cleardoublepage
\phantomsection
\addcontentsline{toc}{chapter}{Daftar Isi}
%\pdfbookmark{\contentsname}{Contents}
\tableofcontents

%%%%%%%%%%%%%%%%%%%%%%%%%%%%%%%%%%%%%%%%%%%%%%%%%
% Bab bab penulisan
%%%%%%%%%%%%%%%%%%%%%%%%%%%%%%%%%%%%%%%%%%%%%%%%
\cleardoublepage
\pagenumbering{arabic}
\chapter{Pendahuluan}
\section{Latar Belakang}
Menulis Latar Belakang wjkajekhahsahd 

\section{Perumusan Masalah}
Berikut rumusan masalah yang ingin saya angkat adalah
\begin{enumerate}
    \item Mengapa ini terjadi?
    \item Bagaimana proses kejadiannya?
    \item Apa saja yang dipengaruhinya?
\end{enumerate}
\section{Tujuan}
Berikut adalah tujuan yang ingin dicapai pada penulisan proposal/TA.
\begin{enumerate}
    \item Untuk mengetahui mengapa ini terjadi;
    \item Untuk mempelajari proses kejadian masalah;
    \item Untuk melihat dampak yang dipengaruhi oleh kejadian ini.
\end{enumerate}
\section{Batassan Masalah}
Hipotesis dari tulisan ini adalah
\begin{enumerate}
    \item Masalah timbul karena A;
    \item Hasil numeriknya menuju $x \rightarrow \infty$
\end{enumerate}

\section{Rencana Kegiatan}
Rencana kegitana yang akan saya lakukan adalah sebagia berikut:
\begin{itemize}
    \item Studi literatur
    \item Memeriksa hasil
\end{itemize}
\section{Jadwal Kegiatan}
The table \ref{Novella} is an example of referenced \LaTeX elements. Laporan proposal ini akan dijadwalkan sesuai dengan tabel yang diberikna berikutnya. 

 
\begin{table}[h!]
  \centering
    \caption{Jadwal kegiatan proposal tugas akhir}
  \label{Novella}
  \begin{tabular}{|c|m{2.5cm}|m{0.01cm}|m{0.01cm}|m{0.01cm}|m{0.01cm}|m{0.01cm}|m{0.01cm}|m{0.01cm}|m{0.01cm}|m{0.01cm}|m{0.01cm}|m{0.01cm}|m{0.01cm}|m{0.01cm}|m{0.01cm}|m{0.01cm}|m{0.01cm}|m{0.01cm}|m{0.01cm}|m{0.01cm}|m{0.01cm}|m{0.01cm}|m{0.01cm}|m{0.01cm}|m{0.01cm}|}
    \hline
    \multirow{2}{*}{\textbf{No}} & \multirow{2}{*}{\textbf{Kegiatan}} & \multicolumn{24}{|c|}{\textbf{Bulan ke-}} \\
    \hhline{~~------------------------}
    {} & {} & \multicolumn{4}{|c|}{\textbf{1}} & \multicolumn{4}{|c|}{\textbf{2}} & \multicolumn{4}{|c|}{\textbf{3}} & \multicolumn{4}{|c|}{\textbf{4}} & \multicolumn{4}{|c|}{\textbf{5}} & \multicolumn{4}{|c|}{\textbf{6}}\\
    \hline
    1 & Studi Literatur & \cellcolor{blue!25} & \cellcolor{blue!25} & \cellcolor{blue!25} & \cellcolor{blue!25}& \cellcolor{blue!25} & \cellcolor{blue!25} & \cellcolor{blue!25} & \cellcolor{blue!25}& \cellcolor{blue!25} & \cellcolor{blue!25} & \cellcolor{blue!25} & \cellcolor{blue!25}& \cellcolor{blue!25} & \cellcolor{blue!25} & \cellcolor{blue!25} & \cellcolor{blue!25}& \cellcolor{blue!25} & \cellcolor{blue!25} & \cellcolor{blue!25} & \cellcolor{blue!25}& \cellcolor{blue!25} & \cellcolor{blue!25} & \cellcolor{blue!25} & \cellcolor{blue!25}\\
    \hline
    2 & Pengumpulan Data & \cellcolor{blue!25} & \cellcolor{blue!25} & \cellcolor{blue!25} & \cellcolor{blue!25} & {} & {} & {} & {} & {} & {} & {} & {}& {} & {} & {} & {}& {} & {} & {} & {}& {} & {} & {} & {}\\
    \hline
    3 & Analisis dan Perancangan Sistem &  {} & {} & {} & {}  & \cellcolor{blue!25} & \cellcolor{blue!25} & \cellcolor{blue!25} & \cellcolor{blue!25} & \cellcolor{blue!25} & \cellcolor{blue!25} & \cellcolor{blue!25} & \cellcolor{blue!25} & {} & {} & {} & {}& {} & {} & {} & {}& {} & {} & {} & {}\\
    \hline
    4 & Implementasi Sistem &  {} & {} & {} & {} & {} & {} & {} & {}& \cellcolor{blue!25} & \cellcolor{blue!25} & \cellcolor{blue!25} & \cellcolor{blue!25} & \cellcolor{blue!25} & \cellcolor{blue!25} & \cellcolor{blue!25} & \cellcolor{blue!25} & {} & {} & {} & {}& {} & {} & {} & {}\\
    \hline
    5 & Analisa Hasil Implementasi &  {} & {} & {} & {} & {} & {} & {} & {}& {} & {} & {} & {} & \cellcolor{blue!25} & \cellcolor{blue!25} & \cellcolor{blue!25} & \cellcolor{blue!25} & \cellcolor{blue!25} & \cellcolor{blue!25} & \cellcolor{blue!25} & \cellcolor{blue!25} & {} & {} & {} & {}\\
    \hline
    6 & Penulisan Laporan & {} & {} & {} & {} & \cellcolor{blue!25} & \cellcolor{blue!25} & \cellcolor{blue!25} & \cellcolor{blue!25}& \cellcolor{blue!25} & \cellcolor{blue!25} & \cellcolor{blue!25} & \cellcolor{blue!25}& \cellcolor{blue!25} & \cellcolor{blue!25} & \cellcolor{blue!25} & \cellcolor{blue!25}& \cellcolor{blue!25} & \cellcolor{blue!25} & \cellcolor{blue!25} & \cellcolor{blue!25}& \cellcolor{blue!25} & \cellcolor{blue!25} & \cellcolor{blue!25} & \cellcolor{blue!25}\\
    \hline
  \end{tabular}

\end{table}


%
\chapter{Kajian Pustaka}

\section{Time Seies method}
Menurut paper Kentang \cite{Kentang}, prsamaan SWE adalah
Berikut diberikan persamaan pengatur dari persamaan gelombang pada gitar

\begin{equation}\label{Pers1}
    a=b+U^{n+1}_{i+1}
\end{equation}
Persamaan (\ref{Pers1}) jadsbahdhavhdvah ajdbajdb
\begin{equation}\label{nama-rumus}
    \int_0^1 \frac{f(x)}{g(x)}\ {\rm dx}=\sin x
\end{equation}

\begin{equation}\label{nama-rumus1}
   \alpha \times \beta =\gamma^{3\alpha}
\end{equation}

\begin{figure}[h!]
    \centering
    \includegraphics[scale=0.1]{Tel-U-Logo.png}
    \caption{Caption}
    \label{fig:my_label1}
\end{figure} 

Rumus (\ref{nama-rumus}) merupakan \textit{contoh} persamaan matematika. persamaan matematika diatas diberi nama \textbackslash label\{nama-rumus\}. dengan $\alpha=\gamma \times 100$

\begin{figure}[h!]
    \centering
    \includegraphics[scale=0.3]{Tel-U-Logo.png}
    \caption{Caption}
    \label{fig:my_label}
\end{figure}

Lihat \textit{pada} Gambar \ref{fig:my_label}

\subsection{Cara memanggil pustaka}
Contoh pustaka prosiding \cite{doyen2014explicit}, jurnal \cite{gunawan2015hydrostatic} dan buku \cite{toro2013riemann}. Atau dapat juga mengguanakan dua pustaka atau lebih dalam \cite{gunawan2015hydrostatic,toro2013riemann}.
%
\chapter{Metodologi dan Desain Sistem}
\section{Flowchart sistem}

%Definisikan bentuk dan warna
\tikzstyle{startstop} = [rectangle, rounded corners, minimum width=3cm, minimum height=1cm,text centered, draw=black, fill=red!30]
\tikzstyle{io} = [trapezium, trapezium left angle=70, trapezium right angle=110, minimum width=3cm, minimum height=1cm, text centered, draw=black, fill=blue!30]
\tikzstyle{process} = [rectangle, minimum width=3cm, minimum height=1cm, text centered, minimum width=3cm, draw=black, fill=orange!30]
\tikzstyle{decision} = [diamond, minimum width=3cm, minimum height=1cm, text centered, draw=black, fill=green!30]
\tikzstyle{arrow} = [thick,->,>=stealth]
\tikzstyle{cloud} = [draw, ellipse,fill=red!20,
    minimum height=2em]

\begin{figure}[h!]
    \centering
    %Mulai menggambar Flowchart
\begin{tikzpicture}[node distance=2cm]
\node (start) [cloud] {Start};
\node (in1) [io, below of=start] {Input};
\node (pro1) [process, below of=in1] {Process 1};
\node (dec1) [decision, below of=pro1] {Decision 1};
\node (pro2a) [process, below of=dec1, yshift=-0.5cm] {Process 2a};
\node (pro2b) [process, right of=dec1, xshift=2cm] {Process 2b};
\node (out1) [io, below of=pro2a] {Output};
\node (stop) [cloud, below of=out1] {Stop};
\draw [arrow] (start) -- (in1);
\draw [arrow] (in1) -- (pro1);
\draw [arrow] (pro1) -- (dec1);
\draw [arrow] (dec1) -- (pro2a);
\draw [arrow] (dec1) -- (pro2b);
\draw [arrow] (pro2b) |- (pro1);
\draw [arrow] (pro2a) -- (out1);
\draw [arrow] (out1) -- (stop);
\end{tikzpicture}
    \caption{Caption flowchart}
    \label{figflow}
\end{figure}

\section{Algoritma}
 Atau dalam bentuk algoritma seperti contoh pada Algoritma \ref{Algo:FVDM} berikut ini:
 

\begin{algorithm}
 \begin{algorithmic}[1]
    \Procedure{FVDM}{$Tfinal, \Delta t$}
    \State \text{Start}
	\State \textbf{For }$n=1:N$\textbf{ do} \Comment{Pemberian nilai awal}
	\State \hspace{0.5cm} Input nilai $x[n]$
	\State \hspace{0.5cm} Input nilai $v[n]$
	\State \textbf{EndFor}
	\State \text{time=0}
	\While{$time < Tfinal$}
	\State \hspace{0.5cm} $time=time +\Delta t$
	\State \hspace{0.5cm} Hitung jarak bamper menggunakan rumus  untuk $n=2,\cdots,N$ 
	\State \hspace{0.5cm} \textbf{If}( $S(n) \leq 0 m)$ \textbf{then return End If}.
	\State \hspace{0.5cm} Tentukan $\lambda$ menggunakan.
	\State \hspace{0.5cm} Hitung kecepatan optimal $v_o(t)$ menggunakan.
	\State \hspace{0.5cm} Hitung percepatan $a_n(time)$ menggunakan .
	\State \hspace{0.5cm} Hitung kecepatan baru dengan $v_n(time)=v_n(time-\Delta t) + a_n(time) \Delta t$.
	\State \hspace{0.5cm} Hitung posisi baru dengan $x_n(time)=x_n(time-\Delta t) + v_n(time) \Delta t$.
	\State \hspace{0.5cm} \textbf{If}( $\Delta v \leq 10^{-5} \&\& a_n(time)\leq 10^{-5})$ \textbf{then}  
    \State \hspace{0.5cm} \hspace{0.5cm} \text{OUTPUT }Cetak hasil data $a_n, v_n, x_n$.
	\State \hspace{0.5cm} \hspace{0.5cm} \textbf{return}.
	\State \hspace{0.5cm} \textbf{End If}.
	\EndWhile
	\State \text{End}
 \EndProcedure
 \end{algorithmic}
 \caption{Prosedur simulasi dinamika lalu lintas menggunakan FVDM.}\label{Algo:FVDM}
\end{algorithm}
%

%%%%%%%%%%%%%%%%%%%%%%%%%%%%%%%%%%%%%%%%%%%%%%%%%%%%%%%%%%%%%%%
% Berikut adalah untuk membuat Daftar pustaka
% Style daftar pustaka dapat disesuaikan dengan mengubah 
%\bibliographystyle{kode}
% dengan kode = acm maka hasinya contoh [1]
% dengan kode = agsm maka hasilnya Harvard style (Sukimin, 2017)
%%%%%%%%%%%%%%%%%%%%%%%%%%%%%%%%%%%%%%%%%%%%%%%%%%%%%%%%%%%%%%%%
\cleardoublepage
\addcontentsline{toc}{chapter}{Daftar Pustaka}
\bibliographystyle{plain} %harvard style
\bibliography{References}
%
%\pagebreak
%%%%%%%%%%%%%%%%%%%%%%%%%%%%%%%%%%%%%%%%%%%%%%%%%%%%%%%%%%%%%%%
% Berikut untuk lampiran
%%%%%%%%%%%%%%%%%%%%%%%%%%%%%%%%%%%%%%%%%%%%%%%%%%%%%%%%%%%%%%
\cleardoublepage
\addcontentsline{toc}{chapter}{Lampiran}
\chapter*{Lampiran}

\end{document}
